\section{AI in the Product: Pricing, Seats and Margins}\label{sec:product}
\IfFileExists{part2/tables/product_macros.tex}{% Generated by code/verify_part2_product.py; do not edit. XBOUNDS placeholders: 2.5, 5, 9.
\providecommand{\prodXbOne}{2.5}
\providecommand{\prodAgentOne}{427}
\providecommand{\prodCopilotOne}{102}
\providecommand{\prodQAOne}{76}
\providecommand{\prodBestOne}{589}
\providecommand{\prodFlowOne}{$+0.65$}
\providecommand{\prodRetOne}{0.75}
\providecommand{\prodAgentPerOne}{17.1}
\providecommand{\prodCopilotPerOne}{8.5}
\providecommand{\prodXbTwo}{5}
\providecommand{\prodAgentTwo}{448}
\providecommand{\prodCopilotTwo}{167}
\providecommand{\prodQATwo}{85}
\providecommand{\prodBestTwo}{685}
\providecommand{\prodFlowTwo}{$+1.40$}
\providecommand{\prodRetTwo}{1.50}
\providecommand{\prodAgentPerTwo}{17.9}
\providecommand{\prodCopilotPerTwo}{14.0}
\providecommand{\prodXbThree}{9}
\providecommand{\prodAgentThree}{482}
\providecommand{\prodCopilotThree}{272}
\providecommand{\prodQAThree}{99}
\providecommand{\prodBestThree}{838}
\providecommand{\prodFlowThree}{$+2.60$}
\providecommand{\prodRetThree}{2.70}
\providecommand{\prodAgentPerThree}{19.3}
\providecommand{\prodCopilotPerThree}{22.7}
\providecommand{\prodVoice}{108}
\providecommand{\prodRankCross}{7.1}
\providecommand{\prodFlowZero}{1/3}
\providecommand{\prodDilOne}{$-0.5$}
\providecommand{\prodDilTwo}{$+2.0$}
\providecommand{\prodDilThree}{$+6.0$}
\providecommand{\prodDilFlip}{3}
}{}% XBOUNDS placeholder values, generated by code/verify_part2_product.py

Sections~\ref{sec:tasks} and \ref{sec:ext} model how the firm produces. This section models what it sells and how AI changes it. The two sides meet in Part~I's parameters: AI products lower gross margin $\gamma$ through inference cost; seat compression lowers retention growth $G_i$; labor budgets that become software budgets raise market growth $G$; better products raise demand pull $\lambda$; and the tariff sets the pass-through $\varphi$ that decides how much of the operational savings of Section~\ref{sec:tasks} the firm keeps. Five conclusions stand out. An agent that does a customer's work is priced as a share of the customer's saving against its cheapest alternative, and under competition from AI-native entrants that price falls to the entrants' cost plus the rent on inputs they lack; data rents fade under diminishing returns (Section~\ref{sec:product-price}). Revenue at risk from seat compression scales with a seat elasticity that no data yet measure; because the loss is permanent, the seat-loss rate at which leaving seat pricing pays is the migration's user cost, which is small when the growth-adjusted discount rate is small, and pure usage is not always the best target (Section~\ref{sec:product-seats}). A point of retention is worth $X-k$ points of gross margin, and falling API prices leave only the rent case for owning a model (Section~\ref{sec:product-portfolio}). AI margins converge to subscription margins only with incomplete pass-through or a scarce input (Section~\ref{sec:product-margin}), and the replacement effect that slows incumbents is small when seat revenue is a small share of the labor it supports (Section~\ref{sec:product-cannibal}). Value pricing and cost-of-revenue automation are complements (Sections~\ref{sec:product-aggregate} and \ref{sec:product-breaks}). Regulation enters the product only as cost and timing, which the examples leave out: an agent sold for a use listed in Annex III of the EU AI Act carries the provider's high-risk obligations from 2 December 2027, and one that is, or is a safety component of, a product under the legislation listed in Annex I from 2 August 2028, the dates set by Regulation (EU) 2026/1744 \citep{eu2024regulation,eu2026regulation}; and the U.S.\ banking agencies' model-risk guidance of April 2026 places generative and agentic AI models outside its scope \citep{fed2026sr262,occ2026bulletin13}, so a bank's validation demands on such an agent come from its own risk management, not from that guidance. Proofs are in Appendix~\ref{app:proofs-product}, and \texttt{code/verify\_part2\_product.py} checks every closed form, threshold and number in the examples.

\subsection{Setting}\label{sec:product-setup}

A unit of a customer's work is an \emph{outcome}: an output unit of an activity or process of Section~\ref{sec:tasks}, such as a resolved ticket or a processed invoice. The vendor sells software that supports this work. It can bill per seat (a person using the product), per unit of usage, per outcome, or with a hybrid tariff, and it can sell an AI agent that does the work itself.

\begin{assumption}[Customer value and adoption]\label{ass:prod-value}
Without the vendor's agent, a customer's cheapest way to produce an outcome costs $\underline c$, the minimum over its feasible modes in house, outsourced and augmented (Proposition~\ref{prop:tasks-modes}(c)). With the agent of the incumbent ($I$) or of an entrant ($E$), the customer bears the residual cost $c^{\rm res}_I=c^Q(e_I)+\eta_I$ or $c^{\rm res}_E=c^Q(e_E)+\eta_E$ per outcome: the quality cost of reviewing the agent's work and bearing its undetected errors at the agent's raw error rate (Proposition~\ref{prop:tasks-review}), plus a handoff and integration cost. The customer's saving per outcome, gross of the price, is $\mathsf V_I:=\underline c-c^{\rm res}_I$ or $\mathsf V_E:=\underline c-c^{\rm res}_E$; without entrants we write $\mathsf V=\mathsf V_I$. Adoption frictions per outcome (integration effort, perceived risk, change cost), as a share of the saving, have a power distribution on $[0,1]$ with exponent $\epsilon_{\rm ad}>0$. At the price $p$ per outcome, the share of in-scope outcomes that moves to the agent is therefore
\begin{equation}\label{eq:prod-D}
D(p)=\Bigl(1-\frac{p}{\mathsf V}\Bigr)^{\epsilon_{\rm ad}}=(1-\kappa_{\rm cap})^{\epsilon_{\rm ad}},\qquad 0\le p\le\mathsf V,
\end{equation}
where $\kappa_{\rm cap}:=p/\mathsf V$ is the vendor's capture share of the saving and $\epsilon_{\rm ad}$ is the elasticity of adoption with respect to the customer's retained share $1-\kappa_{\rm cap}$. The incumbent's full unit cost per outcome is $\hat c_I$: inference $\pi n$, its own escalation and review, hosting and the capital charge. We write $\mathsf c_I:=\hat c_I/\mathsf V_I$.
\end{assumption}

\begin{assumption}[AI-native entrants]\label{ass:prod-entry}
As in Assumption~\ref{ass:ext-entrant}, AI-native entrants offer an agent at full unit cost $\hat c_E$ per outcome and price it at $(1+\upsilon_E)\hat c_E$, where $\upsilon_E\ge 0$ is their markup. Under free entry with replication cost tending to zero, $\upsilon_E\to 0$ (Proposition~\ref{prop:entry}(b)); with economies of scale in foundation models, which can concentrate the market \citep{korinek2025concentrating}, $\upsilon_E>0$. Entrants' agents leave the residual cost $c^{\rm res}_E\ge c^{\rm res}_I$. A customer of the incumbent takes another vendor's agent only if it gains more than a switching premium $p_{\rm sw}$ per outcome. Customers choose the offer with the largest retained saving, and frictions are common to all agents. Define the rent per outcome and the competitive cap
\begin{equation}\label{eq:prod-cap}
\rho_u:=\mathsf V_I-\mathsf V_E+p_{\rm sw},\qquad \bar p:=(1+\upsilon_E)\hat c_E+\rho_u .
\end{equation}
\end{assumption}

Table~\ref{tab:prod-notation} collects the notation. Symbols of Part~I and of Sections~\ref{sec:tasks} and \ref{sec:ext} keep their meanings.

\begin{table}[t]
\centering
\caption{Notation for Section~\ref{sec:product}}\label{tab:prod-notation}
\footnotesize
\begin{tabular}{@{}lp{4.9cm}@{\quad}lp{4.6cm}@{}}
\toprule
Symbol & Meaning & Symbol & Meaning\\
\midrule
$\mathsf V_I,\ \mathsf V_E,\ \underline c$ & customer saving per outcome; best alternative & $\kappa_{\rm cap},\ \epsilon_{\rm ad},\ \mathsf c_I$ & capture share; adoption elasticity; $\hat c_I/\mathsf V_I$\\
$\hat c_I,\ \hat c_E,\ \upsilon_E$ & full unit costs; entrant markup & $\rho_u,\ \bar p$ & rent per outcome; competitive cap\\
$p_{\rm seat},\ \bar p_{\rm seat},\ \gamma_{\rm seat}$ & seat price a year, per unit of work; seat margin & $\mathcal X(t),\ \nu(t),\ \epsilon_{\rm seat}$ & compressed seat share; seat-loss rate; seat elasticity\\
$\mathfrak p,\ \mathfrak q,\ \mathsf B(t)$ & Bass coefficients and curve & $\delta_k,\ e_k,\ \omega_k,\ \varpi_k,\ \nu^*_k$ & tariff $k$: concession, seat exposure, cost, success; threshold\\
$\mathsf f\in\mathbb F;\ \upsilon_{\mathsf f},\ n_{\mathsf f},\ n^E_{\mathsf f}$ & feature; uplift, inference, engineer-years & $\mathsf m;\ \hat c^{\rm net}_{\mathsf m},\ \tilde F_{\mathsf m}$ & model sourcing; net unit cost, fixed cost\\
$\mathsf W,\ \mathsf s$ & labor-to-seat ratio; $\gamma_{\rm seat}/\mathsf W$ & $\nu_E,\ p_L,\ c_{\rm eff}$ & entry hazard; launch price; effective cost\\
$p_{\rm sw};\ \mathsf S_{\rm sw},\ \bar p_{\rm sw},\ \mathrm{cv}_{\rm sw}$ & switching premium; its survival, mean, coefficient of variation & $\bar p_0,\ h_{\rm sw}$ & cap without switching costs; hazard of $p_{\rm sw}$\\
\bottomrule
\end{tabular}
\end{table}

All examples in this section are illustrative. They use a help-desk vendor whose customers' tickets are resolved by customer service representatives (SOC 43-4051). In house a resolution costs \$18.09: the median wage for this occupation in May 2025, \$44{,}770 a year \citep{bls2026oews}, loaded with $\tau_L$, at the 1.7 resolutions an hour of agents without AI assistance in \citet{brynjolfsson2025generative} (Part~II data files). It costs \$12.36 outsourced (at 55\% of in-house cost, with 10\% coordination and a \$0.60 risk premium) and \$15.92 augmented (14\% more resolutions per hour), so $\underline c=\$12.36$. With review time $t^R=0.8t$, a \$30 loss per undetected error, error rates of 4\% for the incumbent's agent and 6.6\% for an entrant's, and integration costs of \$0.05 and \$0.20 a resolution, $\mathsf V_I=\$11.11$ and $\mathsf V_E=\$10.18$. With $p_{\rm sw}=\$0.20$, the rent is $\rho_u=\$1.13$: \$0.78 from data, \$0.15 from the embedded workflow and \$0.20 from switching. Prices are not taken from any one vendor. The examples use ranges of public list prices: \$55 to \$100 per agent a month on annual billing for the middle (professional) seat tiers of four customer-service software products from three companies, and \$0.50 to \$2.00 per resolution or conversation for AI agents.\footnote{\label{fn:prod-prices}List prices, not transaction prices, accessed on 29 October 2026: \citet{freshworks2026pricing}; \citet{fin2026pricing}, for the Intercom helpdesk Advanced plan and the Fin AI Agent; \citet{salesforce2026servicepricing,salesforce2026agentforce}; and \citet{hubspot2026servicepro,hubspot2026customeragent}. Fin, formerly Intercom, has been a Salesforce company since 10 September 2026 \citep{salesforce2026finclose}, so two of the four products share a parent: the ranges are ranges of list prices, not of prices set by independent competitors. The vendors are named only as sources of the ranges; nothing in this section compares or recommends products. Published per-seat tiers run from under \$25 to \$550 a month. The checks use \$25 and \$550 as the ends of that span. The market-conversion and cannibalization conclusions that hold at \$25 also hold at any cheaper seat, because a cheaper seat raises $\mathsf W$ and lowers the seat margin.} Each example reports its conclusion at both ends of the range and says where it changes over the wider span. Values without a named source are assumptions; the seat elasticity and every migration cost are scenario parameters without a measured precedent, and results that use them are conditional on the stated values.

\subsection{Pricing labor-substituting AI}\label{sec:product-price}

\begin{proposition}[Value-based price under competition]\label{prop:prod-price}
Under Assumption~\ref{ass:prod-value}:
\begin{enumerate}
\item[(a)] (Value-based price.) Without entrants, if $\mathsf c_I<1$, the optimal capture share and price are
\begin{equation}\label{eq:prod-kappa}
\kappa^m=\frac{1+\epsilon_{\rm ad}\mathsf c_I}{1+\epsilon_{\rm ad}},\qquad p^m=\frac{\mathsf V_I+\epsilon_{\rm ad}\hat c_I}{1+\epsilon_{\rm ad}} .
\end{equation}
The margin per outcome sold is $(\mathsf V_I-\hat c_I)/(1+\epsilon_{\rm ad})$, profit per in-scope outcome is $\mathsf V_I\,\mathsf C_{\rm ad}(1-\mathsf c_I)^{1+\epsilon_{\rm ad}}$ with $\mathsf C_{\rm ad}:=\epsilon_{\rm ad}^{\epsilon_{\rm ad}}/(1+\epsilon_{\rm ad})^{1+\epsilon_{\rm ad}}$, and if $\hat c_I$ is all variable cost the agent's gross margin is $(1-\mathsf c_I)/(1+\epsilon_{\rm ad}\mathsf c_I)$. If $\mathsf c_I\ge 1$, no price earns a profit.
\item[(b)] (Competition.) Under Assumption~\ref{ass:prod-entry}, the incumbent sells to every adopter at prices up to $\bar p$ and to none above it. If $\bar p>\hat c_I$, its optimal price is $p^*=\min\{p^m,\bar p\}$, so $\kappa^*_{\rm cap}=\min\{\kappa^m,\bar p/\mathsf V_I\}$; if $\bar p\le\hat c_I$, it cannot sell an agent at a profit. Competition binds if and only if $(1+\upsilon_E)\hat c_E<\hat c^{\ddagger}_E:=p^m-\rho_u$.
\item[(c)] (Pass-through.) Where competition does not bind, $dp^*/d\hat c_I=\epsilon_{\rm ad}/(1+\epsilon_{\rm ad})$ and $dp^*/d\mathsf V_I=1/(1+\epsilon_{\rm ad})$, so the elasticity of Lemma~\ref{lem:prices} is $\varphi=\epsilon_{\rm ad}\hat c_I/(\mathsf V_I+\epsilon_{\rm ad}\hat c_I)$. Where it binds, $dp^*/d\hat c_I=0$ and $dp^*/d\hat c_E=1+\upsilon_E$: the incumbent keeps all of its own savings and passes on all of the entrants'.
\item[(d)] (Floor and long-run limit.) A vendor without scarce inputs prices at $(1+\upsilon_E)\hat c_E$. If entry is free ($\upsilon_E\to 0$) and $\hat c_E(t)$ and $\hat c_I(t)$ fall to $c_E^\infty$, the price tends to $\min\{p^m,c_E^\infty+\rho_u\}$ and, where the cap binds, the rent per dollar of revenue to Part~I's $\rho=\rho_u/(c_E^\infty+\rho_u)$; with $\rho_u=0$ the markup vanishes, as in Proposition~\ref{prop:entry}(b).
\item[(e)] (Sources and decay of the rent.) $\rho_u=[c^Q(e_E)-c^Q(e_I)]+(\eta_E-\eta_I)+p_{\rm sw}$: proprietary data act through the error rate, embedded workflows through the handoff cost, and the installed base through the switching premium. Let the error rate fall with the vendor's data stock $\mathcal D$ as $e(\mathcal D)=e_\infty+a_e\mathcal D^{-b_e}$, $a_e,b_e>0$, with review as in Proposition~\ref{prop:tasks-review} (economic share or floor binding). Then the data rent is increasing in $\mathcal D_I$ and at most $L\,a_e(\mathcal D_E^{-b_e}-\mathcal D_I^{-b_e})$, so a constant proportional lead ($\mathcal D_I/\mathcal D_E$ fixed above one) yields a rent that tends to zero as data accumulate. Where the floor binds, $c^Q$ is affine in $e$ and the rent is concave in $\mathcal D_I$.
\item[(f)] (Heterogeneous switching costs.) Let the switching premium vary across customers, independently of frictions, with survival function $\mathsf S_{\rm sw}(x):=\Pr\{p_{\rm sw}\ge x\}$ ($\mathsf S_{\rm sw}=1$ for $x\le 0$), the per-outcome analogue of Proposition~\ref{prop:acq-retention}(a), and let $\bar p_0:=(1+\upsilon_E)\hat c_E+\mathsf V_I-\mathsf V_E$ be the cap without switching costs. At price $p$ the incumbent sells to the share $D(p)\,\mathsf S_{\rm sw}(p-\bar p_0)$; customers with $p_{\rm sw}<p-\bar p_0$ take an entrant's agent if they adopt. An optimal price exists and lies in $[\min\{p^m,\bar p_0\},p^m]$. A single premium (a degenerate distribution) gives (b), with $\bar p=\bar p_0+p_{\rm sw}$. If $p_{\rm sw}$ has a log-concave density on $[0,\infty)$, with hazard $h_{\rm sw}:=-\mathsf S'_{\rm sw}/\mathsf S_{\rm sw}$, and $\hat c_I<\bar p_0<p^m$, the optimal price is $\bar p_0$ if $\ell(\bar p_0)\le h_{\rm sw}(0)$, where $\ell(p):=1/(p-\hat c_I)-\epsilon_{\rm ad}/(\mathsf V_I-p)$, and otherwise the unique root in $(\bar p_0,p^m)$ of
\begin{equation}\label{eq:prod-het}
\ell(p)=h_{\rm sw}(p-\bar p_0).
\end{equation}
At an interior optimum $0<dp^*/d\hat c_I<1$ and $0\le dp^*/d\hat c_E<1+\upsilon_E$: the incumbent passes on part of its own savings and keeps part of the entrants'. With exponential premia of mean $\bar p_{\rm sw}$ the premium earns nothing ($p^*=\bar p_0$) unless $\bar p_{\rm sw}>1/\ell(\bar p_0)$, and at an interior optimum $dp^*/d\hat c_E=0$. With gamma premia of mean $\bar p_{\rm sw}$ and coefficient of variation $\mathrm{cv}_{\rm sw}\le 1$, the optimal price tends to that of (b) with $p_{\rm sw}=\bar p_{\rm sw}$ as $\mathrm{cv}_{\rm sw}\to 0$.
\end{enumerate}
\end{proposition}

\noindent\emph{Intuition.} An agent that does the customer's work is worth the cost of the customer's cheapest alternative, less what the customer still spends on reviewing and integrating it, and the vendor's price is a share of that saving. The share balances margin against adoption: with $\epsilon_{\rm ad}=1$ the vendor splits the surplus over its cost equally with the customer, as with linear demand. This is value-based pricing \citep{hinterhuber2004towards} for an information good with a material marginal cost, to which the bundling results for zero-marginal-cost goods do not carry over \citep{shapiro1999information,bakos1999bundling}. Competition caps the price at an entrant's price plus what the incumbent's agent is worth over the entrant's, a cap that falls one for one with entrants' costs, the defensive logic of Proposition~\ref{prop:ext-should}; what survives free entry is the rent on scarce inputs, as in Proposition~\ref{prop:entry}. The cost-plus floor is a limit, not a forecast: scale economies in foundation models can concentrate supply \citep{korinek2025concentrating,gans2024market}, the leading models' prices extract rent through committed-spend menus while pricing tokens at marginal cost \citep{bergemann2025menu}, and demand for language models shows differentiation and multihoming \citep{fradkin2025demand}; $\upsilon_E>0$ covers these cases. Buyer-side agents that search and contract for customers lower search and switching costs \citep{shahidi2025coasean}, which erodes the switching part of $\rho_u$. Data are nonrival and accumulate with use \citep{jones2020nonrivalry,farboodi2021model,farboodi2023data}, but forecasting gains flatten as data grow \citep{bajari2019impact} and data-savvy entrants can catch up \citep{farboodi2019big,hagiu2023data}; under diminishing returns a fixed proportional data lead is worth less and less, so a data moat must keep widening or be backed by workflows and switching costs. The exception is a market that tips, where the winner's data rent persists \citep{prufer2021competing}. A single switching cost can be charged in full. When switching costs differ across customers, as retention hazards that fall with tenure suggest (Proposition~\ref{prop:acq-retention}), every price above the no-switching cap loses the customers whose costs are lowest, so the price stops where the margin gained equals the hazard of losing them, and some customers switch.

\noindent\emph{Illustrative example.} With $\epsilon_{\rm ad}=1$ and $\hat c_I=\$0.40$, the monopoly capture share is $\kappa^m=0.518$, a price of \$5.75 per resolution at which 48\% of tickets would move to the agent. Entrants at $\hat c_E=\$0.40$ cap the price at $\bar p=\$1.53$, so competition binds: $\kappa^*_{\rm cap}=0.14$, adoption is 86\% and the agent's gross margin is 73.8\%. If entrants' costs fall to a floor of \$0.10, the price tends to \$1.23 and the capture share to 0.11. The competitive price is 8.5\% of the in-house cost per resolution, inside the range of 2.8 to 11.1\% implied by the list prices of \$0.50 to \$2.00 per resolution or conversation (footnote~\ref{fn:prod-prices}); at the two ends of that range the capture share of $\mathsf V_I$ is 0.045 and 0.18, on either side of the model's 0.14. A price set as $\kappa^m$ of the \emph{in-house} saving would be \$8.72, which overstates willingness to pay because outsourcing is cheaper. The data rent of \$0.78 corresponds to a 5.3-fold data lead with inverse-square-root returns ($b_e=1/2$), and it is concave in data here as well; if both firms' data grow tenfold, it falls to \$0.25, and after a hundredfold growth to \$0.08. The single \$0.20 switching premium is the limit $\mathrm{cv}_{\rm sw}\to 0$ of part (f), so these numbers are unchanged. With gamma premia of the same mean $\bar p_{\rm sw}=\$0.20$ and coefficients of variation of 1 (exponential), 0.5 and 0.2 (scenario values; no measured distribution exists), the price is \$1.33, \$1.37 and \$1.45 against \$1.53, the switching part of the rent is \$0, \$0.04 and \$0.12 against \$0.20, and 0, 1.2\% and 0.9\% of the incumbent's customers have premia below its markup over $\bar p_0$, so they would take an entrant's agent rather than its own. With exponential premia the switching premium earns nothing unless its mean exceeds \$1.03. Dispersion lowers the cap, because the customers who would switch most cheaply discipline the price.

\subsection{Seat compression and pricing-model migration}\label{sec:product-seats}

\begin{assumption}[Seats]\label{ass:prod-seats}
A seat-priced line bills $p_{\rm seat}$ per seat a year, or $\bar p_{\rm seat}:=p_{\rm seat}t/t_{\rm FTE}$ per unit of human work, with $t_{\rm FTE}$ paid hours per full-time equivalent. A customer's seats are the full-time equivalents of its staff who use the product: $N_{\rm seat}=\sum_jv_jt^{\rm seat}_{x_j}/t_{\rm FTE}$ over the in-scope activities, with seat-hours per unit $t_j$ in mode $H$, $t^U_j$ in mode $U$, $\xi_jt^R_j$ in mode $A$ and zero in mode $N$, the residual organization of Definition~\ref{def:ext-org}(i), here measured in seat-hours. If customers' operations remove the share $\mathcal X_h(t)$ of in-scope seat-hours at constant volume, seats fall to $N_{\rm seat}(0)[1-\epsilon_{\rm seat}\mathcal X_h(t)]$, where the seat elasticity $\epsilon_{\rm seat}\in[0,1]$ is one when every removed seat-hour removes a seat and lower when customers keep seats for reviewers, supervisors or minimum commitments. No filing reports seats against customers' headcount, so $\epsilon_{\rm seat}$ is a scenario parameter. Let $\mathcal X(t):=\epsilon_{\rm seat}\mathcal X_h(t)$ be the compressed share of seats, $\nu(t):=\dot{\mathcal X}(t)/[1-\mathcal X(t)]$ the proportional seat-loss rate, and $g_v<r$ the growth of the customers' work volume.
\end{assumption}

The Bass curve with innovation coefficient $\mathfrak p>0$ and imitation coefficient $\mathfrak q\ge 0$ \citep{bass1969new} is $\mathsf B(t)=[1-e^{-(\mathfrak p+\mathfrak q)t}]/[1+(\mathfrak q/\mathfrak p)e^{-(\mathfrak p+\mathfrak q)t}]$.

\begin{proposition}[Seat compression]\label{prop:prod-seats}
Under Assumption~\ref{ass:prod-seats}:
\begin{enumerate}
\item[(a)] (Revenue at risk.) Seat revenue is $R_S(t)=p_{\rm seat}N_{\rm seat}(0)e^{g_vt}[1-\mathcal X(t)]$ and grows at $g_v-\nu(t)$: seat compression lowers the line's zero-spend growth by $\nu(t)$. The revenue at risk is $p_{\rm seat}N_{\rm seat}(0)\epsilon_{\rm seat}\mathcal X_h(t)$ a year at constant volume, linear in the seat elasticity, and the seat loss per doubling of the capability horizon is $T_d\nu(t)$ when compression is driven by capability.
\item[(b)] (Customer-side operations.) If customers move activity $j$ from mode $x^0_j$ to $x^1_j$ from the date $T_j$, the later of its feasibility date (Proposition~\ref{prop:ext-feasible}(a)) and its crossover date (Proposition~\ref{prop:tasks-curve}(b)), and adopt along a Bass curve from that date, then $\mathcal X_h(t)=\sum_j\omega^{\rm seat}_j(1-t^{\rm seat}_{x^1_j}/t^{\rm seat}_{x^0_j})\mathsf B(t-T_j)$, where $\omega^{\rm seat}_j$ is activity $j$'s share of seat-hours.
\item[(c)] (Diffusion-limited compression.) If $\mathcal X(t)=\bar x\mathsf B(t)$ with $\bar x=\epsilon_{\rm seat}\bar x_h\in(0,1)$ and $\mathfrak q>\mathfrak p$, seats fall fastest at the tipping date $t^*_B=\ln(\mathfrak q/\mathfrak p)/(\mathfrak p+\mathfrak q)$, when $\mathsf B=(\mathfrak q-\mathfrak p)/(2\mathfrak q)$, at the rate $\bar x(\mathfrak p+\mathfrak q)^2/(4\mathfrak q)$. The seat-loss rate starts at $\nu(0)=\bar x\mathfrak p$ and is single-peaked. If $\mathfrak q>\mathfrak p(1-\bar x)$, it peaks when cumulative adoption reaches
\begin{equation}\label{eq:prod-Bhat}
\hat{\mathsf B}=\frac{1-\sqrt{(1-\bar x)(1+\bar x\mathfrak p/\mathfrak q)}}{\bar x},
\end{equation}
after $t^*_B$, and then falls to zero; otherwise it falls from the start.
\item[(d)] (Capability-limited compression.) If each activity is adopted as soon as it becomes feasible and the doublings it needs are normal with standard deviation $\varsigma_{\rm dbl}$ (Proposition~\ref{prop:ext-feasible}(d)), then $\mathcal X(t)=\bar x\mathsf N((t-t_m)/(\varsigma_{\rm dbl}T_d))$, and seats fall fastest at the median date $t_m$, at the rate $\bar x/(\sqrt{2\pi}\varsigma_{\rm dbl}T_d)$.
\end{enumerate}
\end{proposition}

\noindent\emph{Intuition.} Seat pricing bills for people, so it shrinks with the people that customers' own automation removes, whoever sells the automation. Because the share still exposed shrinks as compression proceeds, the proportional loss rate peaks later than the absolute loss, and it is the proportional rate that drags on retention. Evidence on headcount in exposed occupations is still thin and mixed: early-career employment in the most exposed U.S.\ occupations fell after 2022 \citep{brynjolfsson2025canaries}, while Danish administrative data show no meaningful change in earnings or hours two years after ChatGPT \citep{humlum2025still}.

\noindent\emph{Illustrative example.} Let 40\% of representatives' seat-hours be removable ($\bar x_h=0.4$, an assumption consistent with occupational exposure measures of 0.41 to 0.57 for this occupation; \citealp{tomlinson2025working,eloundou2024gpts}), on \$400 million of seat revenue. Revenue at risk is then $\epsilon_{\rm seat}$ times \$160 million: \$40 million to \$160 million over the scenario range $\epsilon_{\rm seat}\in[0.25,1]$. With the meta-analytic Bass means for consumer durables, $\mathfrak p=0.027$ and $\mathfrak q=0.419$ \citep{vandenbulte2004social}, close to those of the earlier meta-analysis of \citet{sultan1990meta}, and $\epsilon_{\rm seat}=1$, seats fall fastest at 4.7\% a year in year 6.1, and the seat-loss rate rises from 1.1\% a year to a peak of 5.9\% in year 6.7; at $\epsilon_{\rm seat}=0.5$ and $0.25$ the peak is 2.6\% and 1.2\%. Software adoption may be much faster: generative AI reached nearly 40\% of U.S.\ adults within two years of its release \citep{bick2026rapid}, and demand for a new language model levels off within weeks \citep{fradkin2025demand}. In a fast case ($\mathfrak p=0.10$, $\mathfrak q=0.80$, assumed, with $\epsilon_{\rm seat}=1$) the seat-loss rate starts at 4.0\% and peaks at 12.5\% a year in year 2.6.

\begin{assumption}[Pricing-model migration]\label{ass:prod-migrate}
At date $T$ the vendor moves a seat line with gross margin $\gamma$ to a target tariff $k$: pure usage (a price per outcome handled on the platform), a two-part tariff (a platform fee per account plus usage) or a hybrid that keeps a share $e_k\in(0,1)$ of the bill per seat ($e_k=0$ for the first two). Customers accept it at a permanent concession $\delta_k\in[0,1)$ of their current bill, so that
\[
R_k(t)=(1-\delta_k)R_S(T)e^{g_v(t-T)}\Bigl[1-e_k+e_k\frac{1-\mathcal X(t)}{1-\mathcal X(T)}\Bigr],\qquad t\ge T .
\]
The migration costs $\mathcal T_k(T)=\omega_k\gamma R_S(T)$ at $T$ (repricing, billing systems, sales compensation and the transition dip, plus the unwinding of deferred revenue if billing moves to arrears), and succeeds with probability $\varpi_k$; on failure the line reverts to seats. Gross margin is the same under both tariffs. The value of migrating at $T$ is
\[
V_k(T):=\varpi_k\gamma\int_T^\infty e^{-rt}[R_k(t)-R_S(t)]\,dt-e^{-rT}\mathcal T_k(T).
\]
\end{assumption}

\begin{proposition}[Pricing-model migration]\label{prop:prod-migrate}
Under Assumptions~\ref{ass:prod-seats} and \ref{ass:prod-migrate}:
\begin{enumerate}
\item[(a)] $V_k'(T)=e^{-rT}\gamma R_S(T)\bigl\{\varpi_k[\delta_k-(1-\delta_k)(1-e_k)\nu(T)/(r-g_v)]+\omega_k[r-g_v+\nu(T)]\bigr\}$.
\item[(b)] (Seat-compression threshold.) If $\varpi_k(1-\delta_k)(1-e_k)>\omega_k(r-g_v)$, delaying the migration loses value exactly while $\nu(T)>\nu^*_k$, where
\begin{equation}\label{eq:prod-nustar}
\nu^*_k=\frac{(r-g_v)\bigl[\varpi_k\delta_k+\omega_k(r-g_v)\bigr]}{\varpi_k(1-\delta_k)(1-e_k)-\omega_k(r-g_v)} .
\end{equation}
Otherwise $V_k$ is nondecreasing and migrating never pays. The threshold $\nu^*_k$ rises with $\delta_k$, $\omega_k$ and $e_k$, and falls with $\varpi_k$ if $\omega_k>0$.
\item[(c)] (Timing.) If $\nu$ is single-peaked, as in Proposition~\ref{prop:prod-seats}(c), the optimal date is the first date $T_1$ with $\nu(T_1)\ge\nu^*_k$ if $V_k(T_1)>0$, and migration never pays otherwise. In particular $T_1=0$ if $\nu(0)\ge\nu^*_k$.
\item[(d)] (Target tariff.) The best target maximizes $\max_TV_k(T)$ over the menu. This value is nonincreasing in $\delta_k$, $\omega_k$ and $e_k$, and, where positive, nondecreasing in $\varpi_k$, so no target dominates unless it is better in all four.
\end{enumerate}
\end{proposition}

\noindent\emph{Intuition.} Waiting a year costs that year's seat compression, which is permanent and so is capitalized at $1/(r-g_v)$, and saves the user cost of the migration: interest on the one-time cost and on the capitalized concession. Because compression is capitalized while the costs are flows, the threshold is approximately $(r-g_v)\delta_k/[(1-\delta_k)(1-e_k)]$ plus a term of order $(r-g_v)^2$: for any concession and cost, it is small when $r-g_v$ is small, so a firm whose customers are starting to automate should migrate early, phasing contracts at renewal to spread the dip. How early depends on the concession, cost and success parameters, which have no measured precedent. The target need not be pure usage, which protects the whole bill but forfeits the most flat-rate premium: a fixed fee plus usage beats pure usage when metering is costly \citep{sundararajan2004nonlinear}, flat-rate bias raises profit without raising churn while pay-per-use bias raises churn \citep{lambrecht2006paying}, the best access tariff depends on how usage and value are distributed \citep{essegaier2002pricing}, and a metered vendor can lose to a rival that sells outright \citep{balasubramanian2015pricing}. All of this enters $\delta_k$ and $\varpi_k$. Committed-spend contracts, which price the agent's outcomes below the pay-as-you-go rate in exchange for a prepaid budget, are a two-part tariff of this kind \citep{bergemann2025menu}.

\noindent\emph{Illustrative example.} Take $r=0.10$, $g_v=0.04$, a seat margin of 80\% and three targets with scenario parameters: pure usage ($\delta=0.15$, $\omega=0.25$ years of gross profit, $\varpi=0.75$), a two-part tariff ($0.08$, $0.25$, $0.80$) and a hybrid with half the bill per seat ($0.04$, $0.15$, $0.85$). The thresholds are 1.23\%, 0.66\% and 0.65\% a year. With $\epsilon_{\rm seat}=1$ and durable-goods diffusion, the seat-loss rate already exceeds the last two at 1.08\%, so those migrations start now, while pure usage waits 0.32 years. The two-part tariff is then worth the most in both diffusion cases: \$745 million under durable-goods diffusion and \$1{,}033 million under fast diffusion, against \$415 million and \$683 million for pure usage. The conclusion depends on the seat elasticity: at $\epsilon_{\rm seat}=0.5$ the two-part migration waits half a year and is worth \$163 million, and at $0.25$ it never pays (its best date, 2.5 years out, has a value of $-\$110$ million). Its one-time cost is 16\% of the firm's \$500 million revenue if every contract moves at once.

\subsection{The AI feature portfolio and model sourcing}\label{sec:product-portfolio}

A candidate feature $\mathsf f\in\mathbb F$ costs $F_{\mathsf f}$ to develop and uses $n^E_{\mathsf f}$ engineer-years of a capacity $N^E$; both come from the engineering activities of Section~\ref{sec:tasks}, so automating engineering lowers $F_{\mathsf f}$ and raises $N^E$. Among the installed base it is adopted along a Bass curve $\mathsf B_{\mathsf f}(t)$. At full adoption, per dollar of base revenue, it adds a price uplift $\upsilon_{\mathsf f}$, inference cost $\pi(t)n_{\mathsf f}$, retention growth $\delta^G_{\mathsf f}$ and a proportional rise $\delta^\lambda_{\mathsf f}$ in demand pull. Features may need a shared platform (a data layer and an evaluation suite) with cost $F^{\rm plat}$.

\begin{proposition}[The feature portfolio]\label{prop:prod-portfolio}
\begin{enumerate}
\item[(a)] (Value.) To first order, with $X=v_0$, the firm's own weight (Proposition~\ref{prop:ext-customer}(b)), feature $\mathsf f$ adds the score flow
\begin{equation}\label{eq:prod-feature}
R_{\mathsf f}(t)=\mathsf B_{\mathsf f}(t)\bigl[\upsilon_{\mathsf f}-\pi(t)n_{\mathsf f}+(X-k)(\delta^G_{\mathsf f}+\Delta^*\delta^\lambda_{\mathsf f})\bigr]
\end{equation}
per dollar of revenue, and it is worth $\mathrm{NPV}_{\mathsf f}=S\int_0^\infty e^{-(r-g^\circ)t}R_{\mathsf f}(t)\,dt-F_{\mathsf f}$.
\item[(b)] (Portfolio.) Without capacity limits, the firm builds every feature with $\mathrm{NPV}_{\mathsf f}>0$, and builds a shared platform if and only if the positive values of the features that need it sum to more than $F^{\rm plat}$. With capacity $N^E$, the linear relaxation is solved by ranking features by $\mathrm{NPV}_{\mathsf f}/n^E_{\mathsf f}$ and filling capacity in that order, with at most one feature taken fractionally; the value of the integer optimum lies within one feature's NPV of the relaxation. The ratio of the marginal feature is the shadow value of an engineer-year, which prices the automation of engineering in Section~\ref{sec:tasks}.
\item[(c)] (Retention outweighs margin.) A feature whose inference cost exceeds its uplift creates value if and only if, at full adoption, $(X-k)(\delta^G_{\mathsf f}+\Delta^*\delta^\lambda_{\mathsf f})>\pi n_{\mathsf f}-\upsilon_{\mathsf f}$. Ranking features at a screening weight $X^{\rm scr}$ instead of $v_0$ scales the retention part of their value by $(X^{\rm scr}-k)/(v_0-k)$.
\end{enumerate}
\end{proposition}

\noindent\emph{Intuition and illustrative example.} The weight $X-k$ on retention is Part~I's envelope result (Corollary~\ref{cor:priorities}(iii)): the higher the firm's growth weight, the more features that keep customers are worth against features customers pay for. The data do not pin the weight down, so the example is evaluated at illustrative weights of \prodXbOne, \prodXbTwo{} and \prodXbThree. % XBOUNDS: placeholder weights (macros prod*, set in code/verify_part2_product.py)
Take $S=\$500$ million, $k=0$, $\Delta^*=0.05$ and $r-g^\circ=0.05$, a \$15 million platform and 60 engineer-years. An outcome-priced agent (uplift 5\%, inference 1.5\% of base revenue falling at the rate 0.4 a year, retention $+0.1$ point, \$10 million, 25 engineer-years) is worth \$\prodAgentOne, \$\prodAgentTwo{} and \$\prodAgentThree{} million at the three weights; a copilot for human agents that is priced below its inference cost (uplift 0.5\%, inference 0.6\%, retention $+0.2$ point, pull $+2$\%, \$5 million, 12 engineer-years) is worth \$\prodCopilotOne, \$\prodCopilotTwo{} and \$\prodCopilotThree{} million, because its score flow at full adoption is \prodFlowOne, \prodFlowTwo{} and \prodFlowThree{} points, of which \prodRetOne, \prodRetTwo{} and \prodRetThree{} come from retention and pull; the flow is positive for any weight above \prodFlowZero. Quality analytics is worth \$\prodQAOne{} to \$\prodQAThree{} million and a voice agent \$\prodVoice{} million at every weight. At every weight the integer optimum is the agent, the copilot and analytics, which fill 52 engineer-years and are worth \$\prodBestOne, \$\prodBestTwo{} and \$\prodBestThree{} million net of the platform, and the voice agent sets the value of capacity. The ranking by value per engineer-year is what changes: the agent leads at \prodXbOne{} and \prodXbTwo{} (\$\prodAgentPerOne{} and \$\prodAgentPerTwo{} million against the copilot's \$\prodCopilotPerOne{} and \$\prodCopilotPerTwo{} million), and the copilot at \prodXbThree{} (\$\prodCopilotPerThree{} million against \$\prodAgentPerThree{} million); they tie at a weight of \prodRankCross. % XBOUNDS: ranking flips between the second and third placeholder

The agent needs a model. Let $\mathsf m\in\{\mathrm{api},\mathrm{open},\mathrm{own}\}$ be a frontier API, a self-hosted open-weights model fine-tuned on the firm's data, or the firm's own model. At $Q$ outcomes a year, option $\mathsf m$ has variable cost $c^{\rm var}_{\mathsf m}(t)=\pi_{\mathsf m}(t)n_{\mathsf m}+c^h_{\mathsf m}$ per outcome (inference plus orchestration and hosting), capital $K_{\mathsf m}$ per outcome of annual capacity, a fixed cost $\tilde F_{\mathsf m}=F_{\mathsf m}a(r,T_{\mathsf m})+\delta_MF_{\mathsf m}+C^{\rm reval}_{\mathsf m}/T_{\rm up}$ a year (Proposition~\ref{prop:tasks-fullcost}), and rent $\rho_{u,\mathsf m}$. Its net unit cost is $\hat c^{\rm net}_{\mathsf m}:=c^{\rm var}_{\mathsf m}+(r+\delta_K)K_{\mathsf m}-\rho_{u,\mathsf m}$.

\begin{proposition}[Build, buy or open weights]\label{prop:prod-model}
Let the competitive cap bind, so that the price is $(1+\upsilon_E)\hat c_E+\rho_{u,\mathsf m}$.
\begin{enumerate}
\item[(a)] (Volume ladder.) Gross profit is $Q[(1+\upsilon_E)\hat c_E-\hat c^{\rm net}_{\mathsf m}]-\tilde F_{\mathsf m}$, so the optimal option minimizes $Q\hat c^{\rm net}_{\mathsf m}+\tilde F_{\mathsf m}$ and does not depend on entrants' costs. As $Q$ rises, the optimal option moves to lower net unit cost and higher fixed cost, at the thresholds $Q_{\mathsf m\mathsf m'}=(\tilde F_{\mathsf m'}-\tilde F_{\mathsf m})/(\hat c^{\rm net}_{\mathsf m}-\hat c^{\rm net}_{\mathsf m'})$ between neighbors on the lower envelope.
\item[(b)] (Falling API prices.) Let the API's inference price fall at $\delta_\pi$ while the other options are locked to their vintage, their unit costs changing only with a new build (Proposition~\ref{prop:ext-timing}(e)). Then each threshold $Q_{\mathrm{api},\mathsf m}(t)$ is nondecreasing in $t$, and in the long run option $\mathsf m$ is optimal at some volume only if $\rho_{u,\mathsf m}-\rho_{u,\mathrm{api}}>c^{\rm var}_{\mathsf m}+(r+\delta_K)K_{\mathsf m}-c^h_{\mathrm{api}}$: the model must earn a rent that covers its own unit cost.
\item[(c)] (Capital and growth.) Owned capacity raises $k$ by $K_{\mathsf m}Q/S$. By Proposition~\ref{prop:tasks-capital}, value charges it at $r+\delta_K$ and the steady-state cash score at $g+\delta_K$, so for a firm growing faster than $r$ the Rule of 40 penalizes owning compute more than value does.
\item[(d)] (Partner.) A revenue-share partnership that pays the share $\kappa_{\rm par}$ of the price for a partner's model has no fixed cost and unit cost $\kappa_{\rm par}p+c^h$; it is optimal at low volume, and whenever the partner's data or distribution raise $\rho_u$ or $\lambda$ by more than they cost.
\end{enumerate}
\end{proposition}

\noindent\emph{Intuition and illustrative example.} Each option is a line in volume, so the choice is a ladder, as with the modes of Proposition~\ref{prop:tasks-modes}, and run costs shift it toward buying (Proposition~\ref{prop:tasks-buildbuy}); falling API prices then leave only the rent case for owning. Open models currently sell for about 90\% less than comparable closed ones \citep{demirer2025emerging}, which is the unit-cost edge of the middle rung. Let the API cost \$0.40 per outcome, of which \$0.30 is inference, with \$1 million a year of fixed cost; open weights \$0.22 per outcome plus \$0.10 of capital per outcome of annual capacity, \$4 million a year and \$0.10 more rent; and the firm's own model \$0.15 plus \$0.08 of capital, \$15 million a year and \$0.30 more rent ($\delta_K=1/3$). Open weights are cheapest from 12.7 million to 39.5 million outcomes a year, and the own model above that; at 20 million outcomes, open weights win. If the API's inference price falls fivefold a year, open weights lose their unit-cost edge over the API within 0.97 years, and building pays only above 65 million outcomes a year, on the strength of its rent alone.

\subsection{Gross margins of AI products}\label{sec:product-margin}

Let an AI product's unit cost be $u(t)=u_h+u_ce^{-\delta_ut}$ per outcome: a part $u_h$ that does not fall (human escalation, review and support) and a compute part that falls at $\delta_u$, the rate at which the cost of the inference the product actually uses declines. This can be far below $\delta_\pi$ when products move to newer models and use more tokens per task \citep{gundlach2025price}; the preliminary cross-provider elasticity of about $-1.1$ for identical open models \citep{demirer2025emerging} does not show how spending on inference moves with its price, because the authors report that aggregate elasticities are substantially smaller. Fees are the share $f$ of price, and price follows unit cost with the constant pass-through elasticity of Lemma~\ref{lem:prices}, $p(t)=p_0[u(t)/u_0]^{\varphi}$.

\begin{proposition}[Gross-margin path]\label{prop:prod-margin}
\begin{enumerate}
\item[(a)] Gross margin is $\gamma_{\rm AI}(t)=1-f-(1-f-\gamma_0)[u(t)/u_0]^{1-\varphi}$. It rises strictly over time if and only if $\varphi<1$, toward $\gamma_\infty=1-f-(1-f-\gamma_0)(u_h/u_0)^{1-\varphi}$.
\item[(b)] (Crossover.) For a subscription margin $\gamma_{\rm sub}\in(\gamma_0,1-f)$, the AI margin reaches $\gamma_{\rm sub}$ if and only if $\gamma_\infty>\gamma_{\rm sub}$, at
\begin{equation}\label{eq:prod-Tx}
T_\times=\frac1{\delta_u}\ln\frac{u_c}{u_0z^{1/(1-\varphi)}-u_h},\qquad z:=\frac{1-f-\gamma_{\rm sub}}{1-f-\gamma_0};
\end{equation}
with $u_h=0$, $T_\times=\ln(1/z)/[(1-\varphi)\delta_u]$.
\item[(c)] (Limit pricing.) If the competitive cap binds, $\upsilon_E=0$, $\hat c_I=\hat c_E=u$ and $f=0$, then $\gamma_{\rm AI}=\rho_u/(u+\rho_u)$, which rises as $u$ falls, with pass-through elasticity $\varphi=u/(u+\rho_u)$. It converges to subscription margins if and only if $\rho_u/u_h>\gamma_{\rm sub}/(1-\gamma_{\rm sub})$; without a rent it is zero.
\item[(d)] (Blended margin.) If the AI line's revenue share is $w(t)$, the firm's gross margin $w\gamma_{\rm AI}+(1-w)\gamma_{\rm sub}$ falls while $\dot w(\gamma_{\rm sub}-\gamma_{\rm AI})>w\dot\gamma_{\rm AI}$, and at an interior trough the two are equal.
\end{enumerate}
\end{proposition}

\noindent\emph{Intuition and illustrative example.} Falling inference prices raise AI margins only to the extent that prices do not follow them down: with full pass-through the margin is constant, and under competition without a scarce input it is zero. AI product gross margins of 45\% in 2025 and projected values of 53\% in 2026 and 59\% in 2027 \citep{iconiq2026state} identify only the product $(1-\varphi)\delta_u$ if $u_h=0$: a log-linear fit gives 0.147 a year and a crossover with the 74\% median gross margin of Part~I's software sample in 2030. If the cost of the inference used fell fivefold a year ($\delta_u=\ln 5$), that path would imply a pass-through of 0.91; at $\delta_u=0.30$, 0.51. Read instead as constant prices with a cost floor, the same three points give a floor of 23\% of revenue, a limiting margin of 77\% and a crossover in 2033; had the 2027 projection been 57\%, the floor fit would put the limit at 61\%, with no crossover. Three points, two of them projections, cannot separate pass-through from a floor, so the crossover date is fragile. In the help-desk example under limit pricing the margin is 73.8\% now and 91.9\% at the cost floor, because the rent is 11.3 times the floor, against the 2.85 required. If the AI line grows to 35\% of revenue along a Bass curve with $\mathfrak p=0.08$ and $\mathfrak q=0.6$, with the fitted margin path, the firm's blended margin falls from 74\% to a trough of 72.5\% in year 2.7 and recovers by year 5.1.

\subsection{Cannibalization and the replacement effect}\label{sec:product-cannibal}

An incumbent that sells seats earns $\gamma_{\rm seat}\bar p_{\rm seat}$ of gross profit on every unit of human work. If it sells an agent that does the work, it loses that margin on each automated unit. The labor-to-seat ratio $\mathsf W:=\mathsf V_It_{\rm FTE}/(t\,p_{\rm seat})$ measures the saving per seat-year against the seat price, and the replaced margin is the share $\mathsf s:=\gamma_{\rm seat}\bar p_{\rm seat}/\mathsf V_I=\gamma_{\rm seat}/\mathsf W$ of the saving.

\begin{proposition}[Cannibalization]\label{prop:prod-cannibal}
Under Assumption~\ref{ass:prod-value}, with adoption $D$ of \eqref{eq:prod-D}:
\begin{enumerate}
\item[(a)] (Replacement effect.) Let each automated unit release its seat-hours. Without entrants, launching the agent at price $p$ changes the incumbent's profit per in-scope unit by $(p-\hat c_I-\gamma_{\rm seat}\bar p_{\rm seat})D(p)$. The optimal capture share is $\kappa^m_{\rm can}=[1+\epsilon_{\rm ad}(\mathsf c_I+\mathsf s)]/(1+\epsilon_{\rm ad})>\kappa^m$, and launching pays if and only if $\mathsf c_I+\mathsf s<1$, with gain $\mathsf V_I\mathsf C_{\rm ad}(1-\mathsf c_I-\mathsf s)^{1+\epsilon_{\rm ad}}$, below the gain $\mathsf V_I\mathsf C_{\rm ad}(1-\mathsf c_I)^{1+\epsilon_{\rm ad}}$ of an entrant with the same cost.
\item[(b)] (Entry removes it.) Under Assumption~\ref{ass:prod-entry}, if the cap binds, the seats tied to $D(\bar p)$ are lost whether or not the incumbent launches, and launching at $\bar p$ gains $(\bar p-\hat c_I)D(\bar p)$, with no replacement term.
\item[(c)] (Preemption.) Let entry arrive at the Poisson rate $\nu_E$ and work volume grow at $g_v<r$. Before entry, an incumbent that has launched at $p_L$ earns $(p_L-\hat c_I-\gamma_{\rm seat}\bar p_{\rm seat})D(p_L)$ a year more per in-scope unit than one that has not; after entry, a vendor keeps a premium $p_{\rm sw}$ on the customers it served first, so the early launcher earns $p_{\rm sw}D(p_L)$ more. The launch is irreversible. Then launching now is worth $[(p_L-c_{\rm eff})D(p_L)]/(r-g_v+\nu_E)$ more than waiting for entry, with the effective cost
\[
c_{\rm eff}:=\hat c_I+\gamma_{\rm seat}\bar p_{\rm seat}-\frac{\nu_Ep_{\rm sw}}{r-g_v} .
\]
Launching at a later date $T$ before entry is worth $e^{-(r-g_v+\nu_E)T}$ times as much, so the choice is now or at entry. The incumbent should cannibalize now if and only if $\mathsf V_I>c_{\rm eff}$: always if $\mathsf V_I>\hat c_I+\gamma_{\rm seat}\bar p_{\rm seat}$, and otherwise if and only if $\nu_E>\nu_E^*:=(r-g_v)(\hat c_I+\gamma_{\rm seat}\bar p_{\rm seat}-\mathsf V_I)/p_{\rm sw}$. Its launch price is then $p_L=\max\{0,(\mathsf V_I+\epsilon_{\rm ad}c_{\rm eff})/(1+\epsilon_{\rm ad})\}$, below $\kappa^m_{\rm can}\mathsf V_I$.
\end{enumerate}
\end{proposition}

\noindent\emph{Intuition and illustrative example.} A monopolist gains less from an innovation than an entrant because it replaces its own profits \citep{arrow1962economic}. Against the threat of entry, though, it gains more from innovating first, because its monopoly profit exceeds the profit of the duopoly that entry would create \citep{gilbert1982preemptive}; part (c) is a preemption motive of this kind, with entry at a random date and an installed-base premium in place of a patent. Here the replaced profit, the seat margin per unit of work, acts as an extra unit cost; for labor-substituting AI it is small, because a seat costs a small fraction of the labor it supports. At the professional tiers' list prices of \$55 to \$100 per agent a month (footnote~\ref{fn:prod-prices}) and an 80\% margin, the seat margin is \$0.15 to \$0.27 per resolution, 1.3 to 2.4\% of the saving ($\mathsf W=59.5$ to $32.7$). The incumbent's capture share rises only from 0.518 to between 0.525 and 0.530, and the replacement effect costs 2.8 to 5.0\% of the launch gain, so the incumbent should cannibalize at once. The conclusion holds over the whole published span of seat prices, although the cost grows with the seat price: it is 1.3\% at \$25 a month and 25.9\% at \$550, where the seat margin is \$1.49 per resolution, still far below the saving of \$11.11. Entry removes even this, because the seats go either way. The effect binds when seats are expensive relative to the work an agent can take: for a professional seat at \$6{,}000 a year, one hour per unit and an 85\% margin, the seat margin is \$2.45 per unit, and an agent saving \$2.20 at a cost of \$0.50 falls \$0.75 short. With $p_{\rm sw}=\$0.40$ and $r-g_v=0.06$, the incumbent should still launch first if entry is expected within 8.9 years ($\nu_E^*=0.113$); at $\nu_E=0.25$ its effective cost is \$1.29 and it prices at \$1.74 to build a base before entry. The result assumes that the launch is a commitment and that customers do not delay purchases to wait for a better or cheaper agent; when they do, staggered introductions require commitment to future prices and qualities \citep{moorthy1992market}. The model also treats entrants as independent of incumbents and has no acquisitions. In this market, Salesforce agreed on 15 June 2026 to acquire Fin (formerly Intercom), whose AI agent is sold per outcome, for about \$3.6 billion, and completed the acquisition on 10 September 2026 \citep{salesforce2026finagree,salesforce2026finclose}.

\subsection{The product side in Part I's parameters}\label{sec:product-aggregate}

Let the firm have product lines $l$ with revenue $S_l$, revenue shares $w_l=S_l/S$, gross margins $\gamma_l$ and zero-spend growth $G_l$.

\begin{proposition}[Product strategy in Part I's parameters]\label{prop:prod-aggregate}
\begin{enumerate}
\item[(a)] (Mapping.) $\gamma=\sum_lw_l\gamma_l$ and $G_i=\sum_lw_lG_l$ exactly. A seat line has $G_l=g_v-\nu(t)$ plus price growth; a line on a tariff that does not bill per seat has $G_l=g_v$; an outcome-priced line has $G_l=g_v+\dot{\mathsf B}/\mathsf B+d\ln[pD(p)]/dt$. Features raise $G_l$ by $\sum_{\mathsf f}\mathsf B_{\mathsf f}\delta^G_{\mathsf f}$ and demand pull by the proportion $\sum_{\mathsf f}\mathsf B_{\mathsf f}\delta^\lambda_{\mathsf f}$, on top of the operational gain $(1-\delta_\sigma)^{-\alpha}$ of Proposition~\ref{prop:tasks-aggregate}. Moving prepaid seats to billing in arrears raises cash capital intensity by the deferred revenue $k_{\rm def}$ it removes (Remark~\ref{rem:assumptions}(c)).
\item[(b)] (Market.) If customers automate the share $\mathcal X$ of the work and buy the automation at capture share $\kappa_{\rm cap}$, the market for the scope changes from $M_{\rm seat}$ to $M_{\rm seat}[1+\mathcal X(\kappa_{\rm cap}\mathsf W-1)]$. Labor budgets therefore become a larger software market if and only if $\kappa_{\rm cap}>1/\mathsf W$, and market growth $G$ rises by $[\dot{\mathcal X}(\kappa_{\rm cap}\mathsf W-1)+\mathcal X\mathsf W\dot\kappa_{\rm cap}]/[1+\mathcal X(\kappa_{\rm cap}\mathsf W-1)]$; the second term is negative when competition lowers the capture share, the price-deflation channel of Lemma~\ref{lem:prices}(b).
\item[(c)] (Pass-through by tariff.) Seats within the contract term and value-based prices under a binding cap have $\varphi=0$; a monopoly value-based price has $\varphi=\epsilon_{\rm ad}\hat c_I/(\mathsf V_I+\epsilon_{\rm ad}\hat c_I)$; a cost-plus price has $\varphi=1$. A saving $|\Delta u|$ per dollar of revenue in cost of revenue therefore raises $\gamma$ by $(1-\varphi)|\Delta u|$ (Proposition~\ref{prop:tasks-aggregate}(a)).
\item[(d)] (Score.) To first order,
\begin{equation}\label{eq:prod-dR}
\Delta R_X^*=\Delta\gamma-\Delta a-(g^*+\delta_K)\Delta k+(X-k)\Delta G_i+(X-k)\Delta^*\frac{\Delta\lambda}{\lambda},
\end{equation}
exactly linear in $\gamma$ and $G_i$, which enter $R_X^*$ only through $m_0+(X-k)G_i$. A permanent change raises the revenue multiple by $\Delta R^*_{X^\circ}/(r-g^\circ)$ (Proposition~\ref{prop:tasks-aggregate}(c)). Diffusion growth is transient: its time integral is the log of the revenue gain it produces, for the market term $\ln[1+\mathcal X_\infty(\kappa_{\rm cap}\mathsf W-1)]$, so it is bounded, as the score budget of Proposition~\ref{prop:budget} bounds share growth.
\item[(e)] (Complementarity.) Let the firm choose a tariff from a set ordered by decreasing $\varphi$ and a level $z$ of cost-of-revenue automation with saving $|\Delta u(z)|$ and cost $C(z)$. The objective has increasing differences in $(1-\varphi,z)$, so the largest optimal automation level is nondecreasing as the tariff moves toward value-based pricing \citep{milgrom1994monotone}.
\end{enumerate}
\end{proposition}

\begin{corollary}[Margin dilution for growth]\label{cor:prod-dilution}
Let a product plan lower $m_0$ by $D_m>0$ and raise $G_i$ by $\Delta G>0$, with $\lambda$, $k$ and $s$ fixed.
\begin{enumerate}
\item[(i)] $\Delta R_X^*=-D_m+(X-k)\Delta G$ exactly, which is positive if and only if $X-k>D_m/\Delta G$. At small share, the revenue multiple rises to first order if and only if $v_0-k>D_m/\Delta G$ (Proposition~\ref{prop:ext-customer}(b)).
\item[(ii)] The new frontier lies below the old one for $g<\hat g$ and above it for $g>\hat g$, at a unique $\hat g$ on the new domain (or above it everywhere), and $[X(\hat g-\Delta G)-k]\Delta G\le D_m\le[X(\hat g)-k]\Delta G$, where $X(g)$ is the old frontier's dual weight (Proposition~\ref{prop:score}(d)).
\end{enumerate}
\end{corollary}

\noindent\emph{Intuition.} The product side moves mostly growth parameters, weighed at $X-k$, and the operations side mostly cost parameters, weighed at 1. A plan that trades margin for retention raises the score at the firm's own weight and can lower it at an investor's screening weight, so Part~I's two weights disagree exactly on such plans; the frontier rotates, down at low growth and out at high growth. The tariff links the two sides: value-based pricing keeps cost savings and cost-plus pricing passes them on. Market expansion needs only a modest capture share: with a labor saving per seat of 32.7 to 59.5 times the seat price at the professional tiers' list prices, any capture share above 1.7 to 3.1\% enlarges the market.

\noindent\emph{Illustrative example.} With $D_m=0.03$ and $\Delta G=0.01$, the score changes by \prodDilOne, \prodDilTwo{} and \prodDilThree{} points at illustrative weights of \prodXbOne, \prodXbTwo{} and \prodXbThree: the plan pays if and only if $X-k>\prodDilFlip$, so its sign flips between the first two weights and it fails at every weight below \prodDilFlip. % XBOUNDS: placeholder weights; sign flips at X - k = 3
On Part~I's frontier of Figure~\ref{fig:frontier} the new frontier lies below the old for growth below 57\% and above it beyond, where the dual weight exceeds 3.2. In the help-desk example, at $\kappa^*_{\rm cap}=0.14$ and full diffusion, the market for support grows by a factor of 2.2 to 3.5 across the professional seat prices ($\kappa^*_{\rm cap}\mathsf W=4.5$ to $8.2$), and by 6.9 at \$25 a month. The conclusion fails at the top of the published span: at \$550 a month $\kappa^*_{\rm cap}\mathsf W=0.82$ and the market shrinks by 6\%, and conversion needs seats below \$450 a month. It also depends on the outcome price: at the lowest outcome list price ($\kappa_{\rm cap}=0.045$) conversion still holds across the professional tiers ($\kappa_{\rm cap}\mathsf W=1.5$ to $2.7$) but needs seats below \$147, and at the highest ($\kappa_{\rm cap}=0.18$) below \$589. Table~\ref{tab:prod-plan} applies Proposition~\ref{prop:prod-aggregate} to a combined plan in year 3: the two-part tariff now, the agent at the competitive price, the copilot and analytics, and operations savings of 2 points of operating cost and 3 points of cost of revenue. With $\epsilon_{\rm seat}=1$, the plan raises the Rule-of-40 score by 16.7 to 22.7 points under value-based pricing and by 13.7 to 19.7 under cost-plus pricing, at seat prices of \$100 and \$55; the difference of 3 points is the cost-of-revenue saving that cost-plus pricing passes on. A cheaper seat means more work per dollar of seat revenue, so the agent's revenue share is larger (26.0\% against 16.2\%). Most of the gain is transient diffusion growth.

\begin{table}[t]
\centering
\caption{A combined product and operations plan in Part~I's parameters: help-desk example, year 3 (illustrative values)}\label{tab:prod-plan}
\footnotesize
\begin{tabular}{@{}llrrrr@{}}
\toprule
& & \multicolumn{2}{c}{Shift (points)} & \multicolumn{2}{c}{$\Delta R_1^*$ (points)}\\
\cmidrule(lr){3-4}\cmidrule(l){5-6}
Component & Parameter & \$100 & \$55 & \$100 & \$55\\
\midrule
AI agent in the revenue mix (margin 92\%) & $\gamma$ & $+1.89$ & $+3.04$ & $+1.89$ & $+3.04$\\
Features, uplift net of inference & $\gamma$ & $+0.30$ & $+0.27$ & $+0.30$ & $+0.27$\\
Operations, cost of revenue (value pricing; 0 cost-plus) & $\gamma$ & $+3.00$ & $+3.00$ & $+3.00$ & $+3.00$\\
Operations, operating cost & $a$ & $-2.00$ & $-2.00$ & $+2.00$ & $+2.00$\\
Billing in arrears (half the bill), $k_{\rm def}=0.25$ & $k$ & $+8.24$ & $+7.27$ & $-1.24$ & $-1.09$\\
Seat compression avoided by the two-part tariff & $G_i$ & $+2.10$ & $+1.86$ & $+2.10$ & $+1.86$\\
AI expansion within the base (transient) & $G_i$ & $+8.36$ & $+13.43$ & $+8.36$ & $+13.43$\\
Feature retention & $G_i$ & $+0.17$ & $+0.17$ & $+0.17$ & $+0.17$\\
Feature demand pull (proportional) & $\lambda$ & $+1.21$\% & $+1.21$\% & $+0.06$ & $+0.06$\\
\midrule
Total, value pricing / cost-plus & & & & $+16.7$ / $+13.7$ & $+22.7$ / $+19.7$\\
\bottomrule
\end{tabular}

\smallskip
\begin{minipage}{0.95\linewidth}\footnotesize
Notes: columns are the ends of the professional seat tiers' list prices per agent a month (footnote~\ref{fn:prod-prices}); the agent's share of revenue is 16.2\% at \$100 and 26.0\% at \$55. First-order effects \eqref{eq:prod-dR} with $X=1$, $k=0$, $\Delta^*=0.05$ and $g^*=0.15$, at the limit-pricing margin; at the ICONIQ margin path (65\%) the mix row is $-2.47$ at \$100 and $-3.97$ at \$55. At a growth weight $X$ the four growth rows are multiplied by $X$ (by \prodXbOne, \prodXbTwo{} or \prodXbThree{} at the illustrative weights). % XBOUNDS
Durable-goods diffusion of customers' automation, $\epsilon_{\rm seat}=1$ (the compression row scales with $\epsilon_{\rm seat}$), scenario migration parameters; seat revenue \$400 million less the 8\% concession, other lines \$100 million. Values computed by \texttt{code/verify\_part2\_product.py}.
\end{minipage}
\end{table}

\subsection{Breakpoints}\label{sec:product-breaks}

\begin{proposition}[Product breakpoints]\label{prop:prod-breaks}
\begin{enumerate}
\item[(a)] (Competitive price floor.) The price of an undifferentiated agent is $(1+\upsilon_E)\hat c_E$, and competition binds once $(1+\upsilon_E)\hat c_E<\hat c^\ddagger_E$; the agent is viable only if $\hat c_I<\bar p$. If $\upsilon_E=0$ and $\hat c_E(t)=c_E^\infty+[\hat c_E(0)-c_E^\infty]e^{-\delta_\pi t}$ with $c_E^\infty<\hat c^\ddagger_E<\hat c_E(0)$, competition binds from $T_C=\delta_\pi^{-1}\ln\{[\hat c_E(0)-c_E^\infty]/(\hat c^\ddagger_E-c_E^\infty)\}$.
\item[(b)] (Adoption tipping point.) A Bass curve has an interior inflection if and only if $\mathfrak q>\mathfrak p$, at $t^*_B$ and cumulative adoption $(\mathfrak q-\mathfrak p)/(2\mathfrak q)$. Adoption, the growth of AI revenue on a fixed scope and the absolute pace of seat compression peak there; the seat-loss rate peaks later, at \eqref{eq:prod-Bhat}.
\item[(c)] (Other thresholds.) The engine also reports the seat-compression threshold $\nu^*_k$ (Proposition~\ref{prop:prod-migrate}), the margin crossover $T_\times$ (Proposition~\ref{prop:prod-margin}), market conversion $\kappa_{\rm cap}>1/\mathsf W$ (Proposition~\ref{prop:prod-aggregate}(b)), cannibalization $\mathsf V_I>c_{\rm eff}$ (Proposition~\ref{prop:prod-cannibal}), the sourcing volumes $Q_{\mathsf m\mathsf m'}(t)$ (Proposition~\ref{prop:prod-model}) and the dilution weight $D_m/\Delta G$ (Corollary~\ref{cor:prod-dilution}).
\end{enumerate}
\end{proposition}

\noindent Table~\ref{tab:prod-breaks} collects the help-desk values.

\begin{table}[t]
\centering
\caption{Product breakpoints in the help-desk example (illustrative values; $T_C$ for an entrant whose cost starts at \$10.00 and falls fivefold a year to \$0.10; seat results at $\epsilon_{\rm seat}=1$; seat prices \$55 to \$100 a month, with the \$550 top tier where a result changes)}\label{tab:prod-breaks}
\footnotesize
\begin{tabular}{@{}lll@{}}
\toprule
Breakpoint & Threshold & Example\\
\midrule
Competitive price floor & $\hat c_E<\hat c^\ddagger_E=p^m-\rho_u$ & \$4.63; binds now (\$0.40); $T_C=0.49$ years from \$10.00\\
Seat compression & $\nu\ge\nu^*_k$ & 0.66\% (two-part), 1.23\% (usage); $\nu(0)=1.08$\%\\
Margin crossover & $T_\times$ & 2030 (no floor) or 2033 (floor); none if 2027 is 57\%\\
Adoption tipping point & $t^*_B$ & 6.1 years (durables), 2.3 years (fast)\\
Seat-loss peak & $\hat{\mathsf B}$ & 5.9\% a year in year 6.7 (1.2\% at $\epsilon_{\rm seat}=0.25$); 12.5\% (fast)\\
Market conversion & $\kappa_{\rm cap}>1/\mathsf W$ & 1.7 to 3.1\% against $\kappa^*_{\rm cap}=14$\%; fails above \$450 (16.8\% at \$550)\\
Cannibalization & $\mathsf V_I$ vs $\hat c_I+\gamma_{\rm seat}\bar p_{\rm seat}$ & \$11.11 vs \$0.55 to \$0.67 (\$1.89 at \$550): launch now\\
Model sourcing & $Q_{\mathsf m\mathsf m'}$ & 12.7M and 39.5M outcomes; 65M in the long run\\
Dilution & $D_m/\Delta G$ & 3.0: pays at $X-k$ of \prodXbTwo{} and \prodXbThree, not \prodXbOne\\ % XBOUNDS
\bottomrule
\end{tabular}
\end{table}
